\documentclass[11pt]{article}
\usepackage[margin=2.6cm]{geometry}
\usepackage{amsmath,amssymb,amsthm}
\usepackage{enumitem}
\usepackage{xr}
\externaldocument[M-]{model}
\setlist{itemsep=1pt,topsep=3pt}

\newcommand{\C}{\mathcal{C}}
\newcommand{\I}{\mathbb{I}}
\newcommand{\Tm}{\mathrm{Tm}}
\newcommand{\Path}{\mathrm{Path}}
\newcommand{\Fib}{\mathrm{Fib}}
\newcommand{\Cl}{\mathrm{Cl}}
\newcommand{\U}{\mathsf{U}}
\newcommand{\El}{\mathsf{El}}
\newcommand{\Ext}{\mathrm{Ext}}
\newcommand{\horn}{\mathrm{horn}}
\newcommand{\id}{\mathrm{id}}
\newcommand{\D}{\mathsf{D}}
\newcommand{\Lam}{\Lambda}

\theoremstyle{plain}
\newtheorem{theorem}{Theorem}[section]
\newtheorem{lemma}[theorem]{Lemma}
\newtheorem{proposition}[theorem]{Proposition}
\newtheorem{corollary}[theorem]{Corollary}
\theoremstyle{definition}
\newtheorem{definition}[theorem]{Definition}
\theoremstyle{remark}
\newtheorem{remark}[theorem]{Remark}

\title{The localisation of the poset model\\[3pt]
\large a complement to \emph{A model of univalence over a base category with an interval}}
\author{}
\date{17 October 2026}

\begin{document}
\maketitle

\begin{abstract}
\noindent
The poset instance of the model (\texttt{model.tex}, §8.1) is the presheaf
model over finite non-empty posets. Following Sattler, it is localised at an
operation $\D$ built from the inclusion of the linear posets and the strictly
monotone maps (\texttt{\textasciitilde/POSET/poset-model.pdf}). We state the
part of that construction which is proved there, in the setting of
\texttt{model.tex}: $\D$ is a strict monad on types; if $A$ is fibrant then
so is $\D A$; and the two maps $\D A\to\D^2A$ are connected by a path and
are homotopy equivalences. The composition of $\D A$ is the face-uniform
filling of \texttt{model.tex}, Theorem~\ref{M-thm:hornbox}(2), first step, applied to the
restriction of $A$ to the simplex category; no excluded middle is used.
The passage to a left exact modality is the criterion of Coquand, Ruch and
Sattler, which is quoted, not proved.
\end{abstract}

%======================================================================
\section{Linear posets and the restriction to $\Delta$}\label{sec:restr}

We work over the poset instance: $\C$ is the category of finite non-empty
posets, $\Phi(X)$ the decidable sieves, $\Cl=\Fib$. The linear posets
$[n]=\{0<\dots<n\}$ are objects of $\C$, and the simplex category $\Delta$
is the full subcategory on them; $\Delta_+\subseteq\Delta$ has the strictly
monotone maps, which are the composites of face maps. The simplicial
instance of \texttt{model.tex} §8.2 is built on these same maps: its objects
are the numbers $n$, its maps $[m]\to[n]$ the maps of $\C$.

\begin{definition}[restriction]
For a context $\Gamma$ over $\C$, $i^*\Gamma$ is the simplicial set
$n\mapsto\Gamma([n])$; for a type $A$ over $\Gamma$, $i^*A$ is the type over
$i^*\Gamma$ with $(i^*A)(n,\tau)=A([n],\tau)$. Restrictions are those of
$\Gamma$ and $A$ along the maps of $\Delta$.
\end{definition}

A line of $\Gamma$ at $X$ (\texttt{model.tex} §2.2), read at the generalised
points $(f,j):[W]\to X\times[1]$ with $[W]$ linear, gives for every
$f:[n]\to X$ a line of $i^*\Gamma$ at $n$: $(g,j)\mapsto\gamma(fg,j)$. A
partial line of $A$ over $\gamma$ on a sieve $\varphi$ of $X$ gives in the
same way a partial line of $i^*A$ on the sieve $f^*\varphi$ of $[n]$. So a
filling problem for $A$ at $X$ gives a filling problem for $i^*A$ at every
$f:[n]\to X$, and this problem is formed from the data restricted along $f$
only.

%======================================================================
\section{Horn filling from composition}\label{sec:horn}

\begin{lemma}\label{lem:comp-horn}
A fibration structure on $A$ over $\Gamma$ gives a horn filling operation on
$i^*A$ over $i^*\Gamma$ (\texttt{model.tex}, Definition of horn filling),
natural in $\Gamma$.
\end{lemma}

\begin{proof}
Let $\tau\in\Gamma([m+1])$, $0\le k\le m+1$, and a partial element of $i^*A$
on $\Lam^k[m+1]$; it is given by its values $a_Y$ on the faces $Y$ of the
horn, which are linear. If $k=0$, let $q$ be the constant map $0$ and
$g_Y:Y\times[1]\to Y$, $g_Y(y,0)=q(y)$, $g_Y(y,1)=y$. The open box of $A$ at
$[m+1]$ over the constant line at $\tau$, with the sides $a_Y g_Y$ on the
sieve generated by the faces of the horn and the lower lid $a_{\{0\}}\,q$, is
composed from $e_0$ to $e_1$; the end is a filler, since $g_Y(y,1)=y$. If
$k\ge1$, the same with $q(x)=\max(x,k)$, which maps every face of the horn
into itself and $[m+1]$ onto the proper face $\{x\ge k\}$, the roles of $0$
and $1$ exchanged, and the composition from $e_1$ to $e_0$. The construction uses only the
data and the composition, so it is natural.
\end{proof}

This is Theorem~\ref{M-thm:hornbox}(1) of \texttt{model.tex} with the
retraction made explicit, here for a type over posets read on the linear
posets.

%======================================================================
\section{The operation $\D$}\label{sec:D}

\begin{definition}
Let $A$ be a type over $\Gamma$. For $\rho\in\Gamma(X)$, $\D A(X,\rho)$ is the
set of families $u(f)\in A([n],\rho f)$, for all $n$ and $f:[n]\to X$, with
\[
u(f)\,g=u(fg)\qquad\text{for every strictly monotone } g:[m]\to[n].
\]
Restriction along $k:X'\to X$ is $(uk)(f)=u(kf)$. For $\alpha:A\to B$,
$(\D\alpha\,u)(f)=\alpha(u(f))$.
\end{definition}

It suffices to require the condition for the face maps. The family and the
condition are coded (a family indexed by $n$ and the maps $[n]\to X$, a
proposition), so $\D A$ is a type. $\D$ commutes with substitution: for
$s:\Delta\to\Gamma$, $\D(A[s])$ maps to $\D A$ by a cartesian map over $s$
whose fibre maps transport the values along the naturality of $s$, as for
$\Pi$ (\texttt{model.tex}, Theorem~\ref{M-thm:formers}).

\begin{definition}
$\eta_A:A\to\D A$, $(\eta_A a)(f)=af$; and $\mu_A:\D^2A\to\D A$,
$(\mu_A v)(f)=v(f)(\id)$.
\end{definition}

\begin{lemma}[strict monad]\label{lem:monad}
$\eta$ and $\mu$ are natural, $\D\alpha\circ\mu_A=\mu_B\circ\D^2\alpha$, and
\[
\mu_A\circ\eta_{\D A}=\id,\qquad \mu_A\circ\D\eta_A=\id,\qquad
\mu_A\circ\D\mu_A=\mu_A\circ\mu_{\D A},
\]
as equalities of maps, that is of codes.
\end{lemma}

\begin{proof}
Unfold the definitions (poset-model, Lemma~7.4).
\end{proof}

%======================================================================
\section{The composition of $\D A$}\label{sec:comp}

\begin{theorem}\label{thm:DA}
A horn filling operation on $i^*A$ gives a fibration structure on $\D A$.
Taken on the universal family, it is natural in $\Gamma$. With Lemma~\ref{lem:comp-horn}: if $A$ is fibrant, so is
$\D A$.
\end{theorem}

\begin{proof}
Let $\gamma$ be a line of $\Gamma$ at $X$, $u_0\in\D A(\gamma(e))$, and
$(\varphi,u)$ a partial line of $\D A$ ending at $u_0$. For $f:[n]\to X$ the
data read along $f$ (Section~\ref{sec:restr}) form a filling problem for
$i^*A$ at $n$: the line $(g,j)\mapsto\gamma(fg,j)$, the end $u_0(f)$, and the
partial line on $f^*\varphi$ whose value at $(g,j)$ is $u_{fg}$ at the
corresponding chain. These values are natural only along the injective maps
of $\Delta$; they are extended to a simplicial partial line by the
Eilenberg--Zilber factorisation: factor $\langle g,j\rangle=m\sigma$ with
$m$ an injective chain of the prism and $\sigma$ surjective, and take the
value at $m$ restricted along $\sigma$. The factorisation is that of the
composite point of the prism, not of a stage of the partial line: a degenerate
restriction of a stage would otherwise be read through a non-injective map. This is natural along every map,
since the factorisation is unique and $v(m)m'=v(mm')$ for $m'$ injective,
and it commutes with restriction along face maps; only its values on chains
are used below. Let $F(f)$ be its \emph{face-uniform} filler
(\texttt{model.tex}, proof of Theorem~\ref{M-thm:hornbox}(2), first step:
fill the shuffles $P_n,\dots,P_0$ one after the other, each by one horn
filling, and take the side when it contains the identity). The composite is
$f\mapsto F(f)(e')$.
\begin{itemize}
\item It is in $\D A$: a face of the filler at $f$ is the filler at the face
  of $f$, since the face-uniform fillers are uniform along face maps, and
  faces suffice.
\item It agrees with the sides: if $f$ lands in $\varphi$, the side of the
  problem contains the identity and the filler is that side.
\item It is uniform along every $k:X'\to X$: the problem of the restricted
  data at $f$ is the problem of the given data at $kf$, so the two fillers
  are the same.
\end{itemize}
The horn filling operation is only used on the prisms $[n]\times\Delta[1]$,
so it needs no uniformity, and the face-uniform fillers need no excluded
middle.

Naturality is obtained as for $\Ext$ (\texttt{model.tex}, §6.2), by making the
construction once on the universal family. $\El$ is fibrant
(\texttt{model.tex}, §8.1), so $\D\El$ over $\U$ has the fibration just
built. Every fibrant type of the poset model is $c^*\El$ for its code map
$c:\Gamma\to\U$; its $\D$ carries the fibration of $\D\El$ moved along the
cartesian map $\D(c^*\El)\to\D\El$ over $c$ (Section~\ref{sec:D}). For
$s:\Delta\to\Gamma$ the cartesian maps compose, so the structure on
$\D(A[s])$ is the move of that on $\D A$ by the composition law of moves
(\texttt{model.tex}, Lemma~\ref{M-lem:move-laws}); the fillers themselves
are never moved.
\end{proof}

\begin{remark}
Only the uniformity along the face maps is used. This is why $\D A$ is fibrant
without the normalised fillers, which require excluded middle and make
fillers uniform along the degeneracies as well: the condition defining
$\D A$ only concerns $\Delta_+$.
\end{remark}

%======================================================================
\section{Paths and the two maps $\D A\to\D^2A$}\label{sec:paths}

A \emph{homotopy} between maps $\alpha,\beta:A\to B$ over $\Gamma$ is an
element of the path type $\Path_{B}(\alpha,\beta)$ over $\Gamma.A$
(\texttt{model.tex}, Theorem~\ref{M-thm:formers}): at $(\rho,a)$, a line in
$B$ over the constant line at $\rho$ from $\alpha(a)$ to $\beta(a)$, natural in
$(\rho,a)$. We write $\alpha\simeq\beta$; a \emph{homotopy equivalence} is a
map with a two-sided homotopy inverse. Homotopies compose with maps on both
sides.

\begin{lemma}\label{lem:D-htpy}
If $\alpha\simeq\beta$ then $\D\alpha\simeq\D\beta$.
\end{lemma}

\begin{proof}
For a generalised point $F=(f,h):[n]\to X\times[1]$ put
$H'(u)(F)=H(u(f))\langle\id,h\rangle$, where $\langle\id,h\rangle:[n]\to
[n]\times[1]$ (poset-model, Lemma~9.2).
\end{proof}

\begin{lemma}\label{lem:eta-path}
$\eta_{\D A}\simeq\D\eta_A$, by a homotopy which uses no composition or
filling of $A$.
\end{lemma}

\begin{proof}
For $F=(f,h)$, $h$ cuts $[n]=A_0+A_1$; a map $g:[m]\to[n]$ cuts
$[m]=B_0+B_1$, with $g=(g_0+A_1)(B_0+g_1)$. Put
$v(F)(g)=u\bigl(f(g_0+A_1)\bigr)(B_0+g_1)$. Every condition on $v$ is an
instance of one principle: $u(F)K$ is unchanged when $(F,K)$ is
reparametrised along a strictly monotone map (poset-model, Lemma~9.3).
\end{proof}

\begin{corollary}\label{cor:eq}
$\mu_A\circ\D\eta_A=\id$ and $\id\simeq\D\eta_A\circ\mu_A$, with no
fibrancy assumption. If $A$ is fibrant, $\D\eta_A$ and $\eta_{\D A}$ are
homotopy equivalences.
\end{corollary}

\begin{proof}
$\mu_A\circ\D\eta_A=\id$ by Lemma~\ref{lem:monad}. By
Lemma~\ref{lem:eta-path} and Lemma~\ref{lem:D-htpy},
$\D\eta_{\D A}\simeq\D^2\eta_A$; composing with $\mu_{\D A}$ on the left
and using Lemma~\ref{lem:monad},
$\id=\mu_{\D A}\circ\D\eta_{\D A}\simeq\mu_{\D A}\circ\D^2\eta_A
=\D\eta_A\circ\mu_A$. For $\eta_{\D A}$: $\mu_A\circ\eta_{\D A}=\id$, and
$\eta_{\D A}\circ\mu_A\simeq\D\eta_A\circ\mu_A$ by
Lemma~\ref{lem:eta-path}. When $A$ is fibrant so is $\D^2A$
(Theorem~\ref{thm:DA} twice), and homotopies into a fibrant type can be
inverted and composed, by composition in the target; so
$\eta_{\D A}\circ\mu_A\simeq\id$ and $\D\eta_A\circ\mu_A\simeq\id$.
\end{proof}

%======================================================================
\section{The universal property of $\eta$}\label{sec:univprop}

A fibrant type $B$ is \emph{$\D$-modal} when $\eta_B$ is a homotopy
equivalence. By Corollary~\ref{cor:eq}, $\D A$ is $\D$-modal for every
fibrant $A$ (condition 1 of the criterion below).

\begin{proposition}\label{prop:univ}
Let $A$ be a type and $B$ a fibrant $\D$-modal type over $\Gamma$, with $p_B:\D B\to
B$ and $p_B\circ\eta_B\simeq\id$. For $u:A\to B$ put $G(u)=p_B\circ\D u:
\D A\to B$. Then $G(u)\circ\eta_A\simeq u$, and $G(f\circ\eta_A)\simeq f$
for every $f:\D A\to B$. So precomposition with $\eta_A$ is a bijection up to
homotopy between the maps $\D A\to B$ and the maps $A\to B$.
\end{proposition}

\begin{proof}
$G(u)\circ\eta_A=p_B\circ\D u\circ\eta_A=p_B\circ\eta_B\circ u\simeq u$, by
the naturality of $\eta$ (strict). And
$G(f\circ\eta_A)=p_B\circ\D f\circ\D\eta_A\simeq p_B\circ\D f\circ\eta_{\D A}
=p_B\circ\eta_B\circ f\simeq f$, by Lemma~\ref{lem:eta-path} and the
naturality of $\eta$. The homotopies are composed in the fibrant type $B$.
\end{proof}

This is the argument of [CRS], Theorem~1.10, (2)$\Rightarrow$(1), for maps
over $\Gamma$ rather than for the internal function types.

%======================================================================
\section{What is quoted, and what may follow}

With Theorem~\ref{thm:DA}, Lemma~\ref{lem:eta-path}, Corollary~\ref{cor:eq}
and Lemma~\ref{lem:monad}, the criterion of Coquand, Ruch and Sattler
(poset-model, Theorem~10.1) makes $\D$ a left exact modality, and the
$\D$-modal types form a constructive model of univalent type theory with
higher inductive types. That criterion is not formalised. Two further results
of poset-model use only the structures above and could be added: a
proposition with a point over $[0]$ has an element of $\D A$
(Proposition~11.4), and $\D A$ is trivially fibrant when spheres bound in $A$
(Theorem~12.3, the trivial case of Whitehead's principle).

%======================================================================
\appendix
\section{Correspondence with the formalisation}

\noindent Layer L9 of \texttt{seg/core}, over the poset instance; all
statements of this note are formalised except those of the last section.

\medskip
\begin{tabular}{p{4.3cm}p{9.6cm}}
statement & files\\\hline
Section~\ref{sec:restr} & CLocRestr\\
Lemma~\ref{lem:comp-horn} & CLocHorn\\
Definitions of $\D$, $\eta$, $\mu$, Lemma~\ref{lem:monad} & CLocD\\
Theorem~\ref{thm:DA} & CLocEZ (the factorisation), CLocCompV, CLocCompP, CLocComp; naturality on the universal family: CLocDCart, CLocDUniv\\
homotopies, Lemma~\ref{lem:D-htpy} & CHtpyOps (L3: homotopies, whiskering, inverse and composition into a fibrant type), CLocHtpy\\
Lemma~\ref{lem:eta-path} & CLocPath, CLocPathEnd, CLocEtaH\\
Corollary~\ref{cor:eq} & CLocEq\\
Proposition~\ref{prop:univ} & CLocUniv (isModal, dModal; UP.PA.resA, UP.PBm.resB)
\end{tabular}

\end{document}
